% ============================================================
% Αντικατάσταση CAS/Python ενοτήτων — Κεφάλαιο 1
% Γραμμικοί Διανυσματικοί Χώροι
%
% Αυτός ο κώδικας αντικαθιστά τα παλιά casbox + 2×pythonbox
% στο αρχείο ch01_grammikoi_xoroi.tex
%
% QR code: qr_8f4c2704232b.png  (→ latex/qrcodes/)
% URL: https://nikmatz.github.io/ELADIC_GR/code/ch01/
%
% ΣΗΜΕΙΩΣΗ: Δεν χρησιμοποιούμε \begin{ekfonisicas}/\begin{ekfonisipy}
% μέσα στα box — η κάθετη γραμμή τους είναι περιττή εκεί.
% ============================================================

\casactivbox%
  {Maxima: Γραμμικοί Χώροι — Τάξη, Μηδενόχωρος, Gram-Schmidt}%
  {%
    \label{cas:ch1_vector_spaces}
    Δίνεται ο πίνακας
    \[
      A = \begin{pmatrix} 1 & 2 & 1 \\ 2 & 1 & 3 \\ 1 & 3 & -1 \end{pmatrix}.
    \]
    \begin{enumerate}[label=(\alph*),itemsep=2pt]
      \item Υπολογίστε την ορίζουσα $\det(A)$ και την τάξη $\operatorname{rank}(A)$.
        Τι συμπεραίνετε για την αντιστρεψιμότητα του $A$;
      \item Βρείτε την ανηγμένη κλιμακωτή μορφή \texttt{rref(A)}.
        Ποιες στήλες είναι στήλες pivot;
      \item Υπολογίστε τον μηδενόχωρο $\ker(A)$ και τον χώρο στηλών
        $\operatorname{col}(A)$. Επαληθεύστε:
        $\operatorname{rank}(A) + \operatorname{nullity}(A) = 3$.
      \item Εφαρμόστε Gram-Schmidt στα $\mathbf{u}_1=(1,1,0)^\top$,
        $\mathbf{u}_2=(1,0,1)^\top$, $\mathbf{u}_3=(0,1,1)^\top$ και
        επαληθεύστε $\mathbf{e}_i \cdot \mathbf{e}_j = \delta_{ij}$.
    \end{enumerate}
    \smallskip
    \textbf{Εντολές Maxima:}
    \texttt{determinant(A)}, \texttt{rank(A)}, \texttt{rref(A)},
    \texttt{nullspace(A)}, \texttt{columnspace(A)}
  }%
  {https://nikmatz.github.io/ELADIC_GR/code/ch01/}%
  {qr_8f4c2704232b.png}

\subsection*{Δραστηριότητα Python}

\begin{pybox}{Python: Γραμμικοί Χώροι — NumPy, SymPy, Gram-Schmidt}%
             {https://nikmatz.github.io/ELADIC_GR/code/ch01/}%
             {qr_8f4c2704232b.png}
\label{py:ch1_vector_spaces}
Δίνονται τα διανύσματα $\mathbf{v}_1=(1,0,1)^\top$,
$\mathbf{v}_2=(0,1,1)^\top$, $\mathbf{v}_3=(1,1,2)^\top$
και ο πίνακας $A$ της δραστηριότητας Maxima.
\begin{enumerate}[label=(\alph*),itemsep=2pt]
  \item Με \texttt{np.linalg.matrix\_rank} ελέγξτε αν
    $\{\mathbf{v}_1,\mathbf{v}_2,\mathbf{v}_3\}$ είναι γραμμικά ανεξάρτητα.
    Εξηγήστε γιατί η τάξη δεν είναι 3.
  \item Με \texttt{sympy.Matrix.rref()} βρείτε τον μηδενόχωρο του $A$.
    Συγκρίνετε με τα αποτελέσματα Maxima.
  \item Υλοποιήστε Gram-Schmidt για $\mathbf{w}_1=(1,1,0)^\top$,
    $\mathbf{w}_2=(1,0,1)^\top$, $\mathbf{w}_3=(0,1,1)^\top$.
    Κατασκευάστε $Q$ και επαληθεύστε $Q^\top Q = I_3$.
  \item \textit{Προαιρετικό:} Οπτικοποιήστε τα αρχικά και ορθοκανονικά
    διανύσματα σε 3D διάγραμμα (Matplotlib).
\end{enumerate}
\medskip
\textbf{Βιβλιοθήκες / Εντολές:}
\texttt{np.linalg.matrix\_rank()},
\texttt{sympy.Matrix.rref()},
\texttt{sympy.Matrix.nullspace()},
\texttt{np.linalg.norm()},
\texttt{mpl\_toolkits.mplot3d}
\end{pybox}
